\section*{September 7, 2026: Price distributions, property types, and school-site influence}
\textit{Codex draft; diagnostics requested by Jacob.}

The difference between dollar and log estimates survives all 77 single-site
omissions. The descriptive evidence points to differences across more of the
upper half of the price distribution than just the most expensive sales.
Changes in the share of houses versus apartment buildings explain little of
the annual mean-price paths. These findings narrow the explanation; they do
not establish that closure caused the price changes.

The analysis keeps equal weight per transaction and the quarter-mile comparison
of 30 closed and 49 still-open schools. The 10,921 sales cover 29 treated and
48 control sites during 2008--2018. Each omission removes the transactions
assigned to one physical school site without reassigning them. The eight
annual models combine dollars or logs, pooled DiD or event study, and the
presence or absence of property controls. The full-sample estimates and
clustered standard errors reproduce the baseline before any sites are omitted.
No additional prices are trimmed.

\textbf{Means and percentiles.} As an illustrative interval, between 2012 and
2016 the treated mean rose from \$238,759 to \$351,113, while the control mean
rose from \$149,427 to \$198,729. This widens the mean gap by \$63,051.
Median prices rose by almost the same dollar amount: \$45,199 near closed
schools and \$47,177 near controls. The treated 75th percentile rose by
\$185,888 versus \$48,879 in controls. These are unadjusted distributional
comparisons, not estimates of causal effects at individual quantiles.
The plots also show differences in the groups' recovery before 2013.

\textbf{Expensive sales.} The annual top 5\% account for 18.0--23.4\% of
sale value near treated schools and 21.2--26.5\% near controls. Between 2012
and 2016 their contribution to the widening mean gap is \$11,258 of
\$63,051; the remaining contribution comes from sales below that annual
cutoff. Thus a few extreme transactions alone do not explain that descriptive
change. Of the 170 treated and 387 control sales above their group-year
95th percentile, 54.7\% and 34.9\%, respectively, are two-to-six-unit
apartment buildings. The saved transaction list includes class codes,
building sizes, apartment counts, and school assignments for inspection.

\textbf{School-site influence.} With property controls, the pooled dollar
DiD is \$56,983. Omitting Humboldt/Duprey lowers it to \$39,555;
omitting Lafayette gives \$44,435, Peabody \$45,379, and Trumbull \$49,078.
Across all omissions it ranges from \$39,555 to \$64,552. The corresponding
log estimate ranges from 0.0167 to 0.0856, compared with 0.0552 in the full
sample. Estimates remain imprecise: the full-sample dollar interval and those
omitting each of the four sites above include zero. Positive point estimates
under every single-site omission do not rule out concentration across several
sites together. The event-study figure shows that the contrasting annual
log and dollar paths also persist.

\textbf{Property mix.} Apartment buildings comprise 57.0\% of treated sales
before closure (2008--2012) and 53.9\% afterward (2014--2018), versus 47.5\%
and 43.5\% for controls. Holding each group's two property-type shares at
its pre-closure shares leaves its annual means nearly unchanged, except for
a \$13,496 upward adjustment to the treated mean in 2018. Property size,
quality, location, and selection into sale can still change within each type.

{\small Reproduction: run Make in
\texttt{tasks/audits/home\_price\_distribution/code}. Seven saved datasets have
automatic reports; five figures follow. The audit uses the existing local
source vintage. Sample counts and mean prices match the baseline in every
site-year. Independent dummy-variable OLS checks reproduce pooled coefficients
for the full sample and the Humboldt/Duprey omission.}

\clearpage
\begin{center}\includegraphics[width=\textwidth]{../input/annual_price_distribution.png}\end{center}
The same dollar scales apply to treated and control panels within each row.
The median and lower percentiles have different paths from the upper
percentiles; annual means need not summarize a common shift in all prices.

\clearpage
\begin{center}\includegraphics[width=\textwidth]{../input/expensive_sales_contribution.png}\end{center}
The two shaded contributions sum to each annual mean. Both use the total
number of sales as denominator. Upper-tail membership is defined separately
within each treatment group and year, so these are changing sets of properties.

\clearpage
\begin{center}\includegraphics[width=.95\textwidth,height=.85\textheight,keepaspectratio]{../input/property_composition.png}\end{center}
The fixed-share comparison adjusts only the house/apartment distinction.
It does not hold the mix of locations or characteristics within each type fixed.

\clearpage
\begin{center}\includegraphics[width=\textwidth]{../input/site_omission_did.png}\end{center}
Each panel displays its eight most influential omissions, selected by the
absolute change in the point estimate. The data files contain all 77 omissions.
Humboldt and Duprey share a site and are omitted together. Intervals use
standard errors clustered by the remaining school sites.

\clearpage
\begin{center}\includegraphics[width=\textwidth]{../input/site_omission_events.png}\end{center}
The shaded range describes sensitivity of point estimates to single-site
omissions. It is not a confidence band. Vertical intervals apply to the
full-sample estimates, and 2012 is the omitted reference year.
